% =====================================================================
% Relatorio do Marco 2 (OCRA) -- MATA60 -- Caso ScienceReview
% Formatacao segundo "Instructions for Authors of SBC Book Chapters" (Regra 4).
% Secoes seguem o Quadro 3 (secao 4.4 do enunciado).
% Comentarios "BAREMA" indicam o que cada secao precisa ter para pontuar
% na rubrica 4.5; comentarios "REGRA" indicam penalidades das regras 4.1.
% =====================================================================
%
% RUBRICA 4.5 (nota 0-10):
%   P1. Readequacao e Normalizacao (peso 3)
%     Q1: plano de readequacao e normalizacao, a partir do feedback da
%         Entrega 1, fundamentado e associado sistematicamente aos requisitos
%     Q2: Q1 + plano seguido, gerando novos modelos conceitual e logico
%     Q3: Q2 + novo banco corrigido e populado conforme o plano
%   P2. Consultas (peso 4)
%     Q1: >= 30 consultas intermediarias, todas ligadas a RF com justificativa
%     Q2: Q1 + >= 20 consultas avancadas, todas ligadas a RF com justificativa
%   P3. Otimizacao (peso 3)
%     Q1: 20 consultas escolhidas (10 intermediarias + 10 avancadas),
%         descritas quanto a complexidade, com plano de otimizacao
%         (algebra relacional e indexacao) para cada uma
%     Q2: Q1 + as duas solucoes aplicadas + cada consulta executada >= 20
%         vezes por plano, usando as consultas sem alteracao como baseline
%     Q3: Q2 + relatorio compara cada plano com o baseline usando
%         speedup medio e desvio padrao
%
% REGRAS 4.1:
%   R1 diagramas conceituais no BrModelo, notacao Chen  -> invalida a entrega
%   R2 anexos/links publicos, sem autenticacao           -> invalida a entrega
%   R3 imagens legiveis                                  -> invalida a entrega
%   R4 modelo SBC de capitulo de livro                   -> -3 pontos
%   R5 reproducao so com PostgreSQL + pgAdmin            -> -3 pontos
%   R6 SQL ISO/IEC 9075, sem Python/R/Java etc.          -> -3 pontos
%   R7 intermediaria: >= 3 tabelas e >= 2 de JOIN, GROUP BY, WINDOW, COUNT
%                                                        -> -0,25 por consulta
%   R8 avancada: >= 3 tabelas e >= 3 de SUBCONSULTA, JOIN, GROUP BY,
%      WINDOW, COUNT                                     -> -0,25 por consulta
%   R9 cada consulta associada a pelo menos um RF, explicito no relatorio
%                                                        -> -0,25 por consulta
% =====================================================================
\documentclass[12pt,a4paper]{article}

\usepackage[T1]{fontenc}
\usepackage[utf8]{inputenc}
\usepackage[brazil]{babel}
\usepackage{mathptmx}
\usepackage[scaled=0.92]{helvet}
\usepackage[a4paper,top=3.5cm,bottom=2.5cm,left=3cm,right=3cm]{geometry}
\usepackage{graphicx}
\usepackage{titlesec}
\usepackage{caption}
\usepackage{booktabs}
\usepackage{longtable}
\usepackage{array}
\usepackage{url}
\usepackage{amsmath}
\usepackage{listings}
\usepackage[square]{natbib}
\usepackage[hidelinks]{hyperref}

% ---------- Pagina e paragrafos ----------
\pagestyle{empty}
\setlength{\parindent}{1.27cm}
\setlength{\parskip}{6pt}

% ---------- Secoes ----------
\titleformat{\section}{\normalfont\fontsize{13}{15}\selectfont\bfseries}{\thesection.}{0.5em}{}
\titlespacing*{\section}{0pt}{18pt}{6pt}
\titleformat{\subsection}{\normalfont\fontsize{12}{14}\selectfont\bfseries}{\thesubsection.}{0.5em}{}
\titlespacing*{\subsection}{0pt}{12pt}{6pt}

% ---------- Figuras, tabelas e legendas ----------
\captionsetup{font={footnotesize,sf,bf},labelsep=period,justification=justified,
  singlelinecheck=true,margin=0.8cm,skip=6pt}
\captionsetup[table]{position=top}

% ---------- Citacoes no estilo [Autor Ano] ----------
\setcitestyle{square,aysep={},yysep={,}}

% ---------- Codigo SQL ----------
\lstset{language=SQL,basicstyle=\ttfamily\footnotesize,breaklines=true,
  columns=fullflexible,keepspaces=true,frame=single,showstringspaces=false,
  upquote=true}

% ---------- Utilitarios ----------
\newcommand{\pendente}[1]{\textbf{[PENDENTE: #1]}}
\newcommand{\tb}[1]{\texttt{#1}}
\newcommand{\rel}[1]{\textsf{#1}}

\begin{document}
\thispagestyle{empty}

% =====================================================================
% Primeira pagina
% =====================================================================
\begin{flushleft}
{\fontsize{20}{24}\selectfont\bfseries ScienceReview: readequação, normalização e otimização de consultas no banco de dados de gestão de submissões e avaliação científica}\\[1cm]
{\large Filipe dos Santos Freitas, Larissa Vitória Maturino Assis, Leticia Silva Oliveira, Lucas dos Santos Lima e Lucy Cassia Cerqueira Silva}\\[4pt]
{\normalsize Instituto de Computação -- Universidade Federal da Bahia (UFBA)\\
Av. Milton Santos, s/n -- Campus de Ondina -- Salvador -- BA, 40170-110}
\end{flushleft}

\vspace{1cm}

\begin{center}\textbf{\textit{Abstract}}\end{center}
\vspace{-6pt}
{\leftskip=0.8cm\rightskip=0.8cm\noindent\itshape
This chapter presents the second iteration of the ScienceReview database. Based on
the feedback from the first delivery, the conceptual and logical models were revised
and normalized, and a new database was deployed and populated in PostgreSQL. Fifty
SQL queries (30 intermediate and 20 advanced), each linked to at least one
requirement, were written for the new database. Twenty of them were submitted to two
optimization strategies, relational algebra rewriting and indexing, and compared
against the unmodified queries using mean speedup and standard deviation over
repeated executions.\par}

\begin{center}\textbf{\textit{Resumo}}\end{center}
\vspace{-6pt}
{\leftskip=0.8cm\rightskip=0.8cm\noindent\itshape
Este capítulo apresenta a segunda iteração do banco de dados do ScienceReview. A
partir do retorno da primeira entrega, os modelos conceitual e lógico foram
revisados e normalizados, e um novo banco foi implantado e populado em PostgreSQL.
Foram escritas 50 consultas SQL (30 intermediárias e 20 avançadas), cada uma
associada a pelo menos um requisito. Vinte delas foram submetidas a duas estratégias
de otimização, reescrita por álgebra relacional e indexação, e comparadas às
consultas sem alteração por meio do \textit{speedup} médio e do desvio padrão em
execuções repetidas.\par}

% =====================================================================
\section{Melhorias e normalização}\label{sec:melhorias}
% =====================================================================
% BAREMA P1-Q1: plano de readequacao e normalizacao, a partir do feedback
% da Entrega 1, fundamentado e associado sistematicamente aos requisitos.

\subsection{Retorno da primeira entrega}\label{sec:feedback}
% Cada ponto do feedback -> acao tomada -> RF afetados.

{\small
\setlength{\tabcolsep}{4pt}
\begin{longtable}{>{\raggedright\arraybackslash}p{0.7cm}
                  >{\raggedright\arraybackslash}p{4.6cm}
                  >{\raggedright\arraybackslash}p{6.0cm}
                  >{\raggedright\arraybackslash}p{2.0cm}}
\caption{Retorno da Entrega 1 e ações tomadas.}\label{tab:feedback}\\
\toprule
\textbf{\#} & \textbf{Ponto apontado} & \textbf{Ação} & \textbf{RF}\\
\midrule
\endfirsthead
\toprule
\textbf{\#} & \textbf{Ponto apontado} & \textbf{Ação} & \textbf{RF}\\
\midrule
\endhead
\bottomrule
\endfoot
 & & & \\
\end{longtable}
}

\subsection{Plano de readequação}\label{sec:plano}
% Problemas do modelo da Entrega 1 (do feedback e da revisao do grupo),
% correcao planejada, RF associados e forma normal envolvida.

{\small
\setlength{\tabcolsep}{4pt}
\begin{longtable}{>{\raggedright\arraybackslash}p{0.7cm}
                  >{\raggedright\arraybackslash}p{4.0cm}
                  >{\raggedright\arraybackslash}p{5.4cm}
                  >{\raggedright\arraybackslash}p{1.6cm}
                  >{\raggedright\arraybackslash}p{1.4cm}}
\caption{Plano de readequação.}\label{tab:readequacao}\\
\toprule
\textbf{\#} & \textbf{Problema} & \textbf{Correção} & \textbf{RF} & \textbf{FN}\\
\midrule
\endfirsthead
\toprule
\textbf{\#} & \textbf{Problema} & \textbf{Correção} & \textbf{RF} & \textbf{FN}\\
\midrule
\endhead
\bottomrule
\endfoot
 & & & & \\
\end{longtable}
}

\subsection{Normalização}\label{sec:normalizacao}
% Aplicacao e justificativa das estrategias de normalizacao (Quadro 3):
% dependencias funcionais de cada relacao e verificacao de 1FN, 2FN, 3FN
% e FNBC. Desnormalizacoes mantidas precisam de justificativa.

% =====================================================================
\section{Novo banco de dados}\label{sec:novobd}
% =====================================================================
% BAREMA P1-Q2: plano seguido, gerando novos modelos conceitual e logico.
% BAREMA P1-Q3: novo banco corrigido e populado conforme o plano.

\subsection{Modelo conceitual}
% REGRA R1: diagrama feito no BrModelo, notacao de Peter Chen.
% REGRA R3: imagem legivel.

\subsection{Modelo lógico}

\subsection{Implantação}
% Script DDL e estrategia de traducao do MER para o relacional.
% Manter pelo menos duas tabelas com prefixo "tabela_imp" (regra do Marco 1).

\subsection{População}
% Volume por tabela. Entrega 1 exigia >= 3 tabelas com >= 5.000 registros.

% =====================================================================
\section{Consultas}\label{sec:consultas}
% =====================================================================
% BAREMA P2-Q1: >= 30 intermediarias ligadas a RF, com justificativa.
% BAREMA P2-Q2: >= 20 avancadas ligadas a RF, com justificativa.
% REGRA R7, R8: conferir tabelas e recursos de cada consulta.
% REGRA R9: RF explicito para cada consulta.
% REGRA R6: SQL padrao ISO/IEC 9075.
% Codigo comentado das consultas: aqui ou no Anexo C.

{\small
\setlength{\tabcolsep}{3pt}
\begin{longtable}{>{\raggedright\arraybackslash}p{0.8cm}
                  >{\raggedright\arraybackslash}p{5.4cm}
                  >{\raggedright\arraybackslash}p{0.8cm}
                  >{\raggedright\arraybackslash}p{1.8cm}
                  >{\raggedright\arraybackslash}p{1.6cm}
                  >{\raggedright\arraybackslash}p{2.4cm}}
\caption{Catálogo das consultas.}\label{tab:consultas}\\
\toprule
\textbf{ID} & \textbf{Objetivo} & \textbf{Tab.} & \textbf{Recursos} & \textbf{RF} & \textbf{Justificativa}\\
\midrule
\endfirsthead
\toprule
\textbf{ID} & \textbf{Objetivo} & \textbf{Tab.} & \textbf{Recursos} & \textbf{RF} & \textbf{Justificativa}\\
\midrule
\endhead
\bottomrule
\endfoot
\multicolumn{6}{l}{\textit{Consultas intermediárias}}\\
 & & & & & \\
\midrule
\multicolumn{6}{l}{\textit{Consultas avançadas}}\\
 & & & & & \\
\end{longtable}
}

% =====================================================================
\section{Plano de otimização 1: álgebra relacional}\label{sec:otim-algebra}
% =====================================================================
% BAREMA P3-Q1: 20 consultas escolhidas (10 intermediarias + 10 avancadas),
% criterio de escolha justificado, complexidade descrita e plano de
% otimizacao para cada consulta.
% Enunciado: pelo menos 3 tecnicas diferentes de algebra relacional.

\subsection{Seleção das consultas}
% Criterios sugeridos pelo enunciado: (i) cumprem mais requisitos;
% (ii) ligadas a requisitos nucleares; (iii) maior custo no EXPLAIN ANALYZE.

{\small
\setlength{\tabcolsep}{4pt}
\begin{longtable}{>{\raggedright\arraybackslash}p{0.9cm}
                  >{\raggedright\arraybackslash}p{4.6cm}
                  >{\raggedright\arraybackslash}p{4.4cm}
                  >{\raggedright\arraybackslash}p{3.2cm}}
\caption{Consultas escolhidas para otimização.}\label{tab:selecao}\\
\toprule
\textbf{ID} & \textbf{Critério de escolha} & \textbf{Complexidade} & \textbf{Custo original}\\
\midrule
\endfirsthead
\toprule
\textbf{ID} & \textbf{Critério de escolha} & \textbf{Complexidade} & \textbf{Custo original}\\
\midrule
\endhead
\bottomrule
\endfoot
 & & & \\
\end{longtable}
}

\subsection{Técnicas de álgebra relacional}

\subsection{Aplicação por consulta}

% =====================================================================
\section{Plano de otimização 2: indexação}\label{sec:otim-indice}
% =====================================================================
% BAREMA P3-Q1: plano de indexacao para cada consulta escolhida.
% BAREMA P3-Q2: plano aplicado.

\subsection{Estratégia de indexação}

\subsection{Aplicação por consulta}

% =====================================================================
\section{Análise de desempenho}\label{sec:desempenho}
% =====================================================================
% BAREMA P3-Q2: cada consulta executada >= 20 vezes em cada plano;
% baseline = consultas sem alteracao.
% BAREMA P3-Q3: comparacao de cada plano com o baseline usando speedup
% medio e desvio padrao.

\subsection{Metodologia}
% Ambiente (hardware, versao do PostgreSQL), numero de execucoes,
% forma de medir o tempo e definicao do speedup.

\subsection{Resultados}

{\small
\setlength{\tabcolsep}{4pt}
\begin{longtable}{>{\raggedright\arraybackslash}p{0.9cm}
                  >{\raggedleft\arraybackslash}p{2.0cm}
                  >{\raggedleft\arraybackslash}p{2.0cm}
                  >{\raggedleft\arraybackslash}p{2.0cm}
                  >{\raggedleft\arraybackslash}p{2.6cm}
                  >{\raggedleft\arraybackslash}p{2.6cm}}
\caption{Desempenho de cada plano em relação ao baseline.}\label{tab:resultados}\\
\toprule
\textbf{ID} & \textbf{Baseline (ms)} & \textbf{Álgebra (ms)} & \textbf{Índices (ms)} & \textbf{Speedup álgebra ($\bar{S}\pm\sigma$)} & \textbf{Speedup índices ($\bar{S}\pm\sigma$)}\\
\midrule
\endfirsthead
\toprule
\textbf{ID} & \textbf{Baseline (ms)} & \textbf{Álgebra (ms)} & \textbf{Índices (ms)} & \textbf{Speedup álgebra ($\bar{S}\pm\sigma$)} & \textbf{Speedup índices ($\bar{S}\pm\sigma$)}\\
\midrule
\endhead
\bottomrule
\endfoot
 & & & & & \\
\end{longtable}
}

\subsection{Discussão}

% =====================================================================
% Referencias
% =====================================================================
\renewcommand{\refname}{Referências}
\bibliographystyle{apalike}
\bibliography{referencias_marco2}

% =====================================================================
% Anexos
% =====================================================================
% REGRA R2: links externos publicos, sem autenticacao.
% REGRA R5: reproducao apenas com PostgreSQL + pgAdmin.
\section*{Anexo A -- Script DDL}

\section*{Anexo B -- Script de população}

\section*{Anexo C -- Consultas}

\section*{Anexo D -- Scripts de otimização e medição}

\section*{Anexo E -- Roteiro de reprodução}

\section*{Anexo F -- Materiais externos}

\end{document}
